%==============================================================================
% Makalah Gemastik — Buletin Pagelaran Mahasiswa Nasional Bidang TIK
% Judul: Peningkatan Keandalan dan Efisiensi Token LLM pada Data Time Series
%        melalui Pendekatan Verifikasi Deterministik
% Kompilasi: pdflatex 02_KrocoMasStanis.tex  (jalankan 2x untuk referensi)
% Butuh berkas gambar: arsitektur_sistem.png (di folder yang sama)
%==============================================================================
\documentclass[10pt,a4paper,twocolumn]{article}

\usepackage[T1]{fontenc}
\usepackage[utf8]{inputenc}
\usepackage[indonesian]{babel}
\usepackage{mathptmx}            % Times New Roman untuk teks & matematika
\usepackage{amsmath,amssymb}
\usepackage{graphicx}
\usepackage{booktabs}
\usepackage{array}
\usepackage{ragged2e}
\usepackage{caption}
\usepackage{titlesec}
\usepackage{fancyhdr}
\usepackage{url}
\usepackage[hidelinks]{hyperref}
\usepackage[a4paper,left=1.8cm,right=1.8cm,top=2.6cm,bottom=2.2cm,%
            columnsep=0.7cm,headsep=0.5cm]{geometry}

%------------------------------------------------------------------------------
% Penomoran bagian: I., A. (IEEE style)
%------------------------------------------------------------------------------
\renewcommand{\thesection}{\Roman{section}}
\renewcommand{\thesubsection}{\Alph{subsection}}

\titleformat{\section}
  {\normalfont\small\bfseries}{\thesection.}{0.5em}{\MakeUppercase}
\titlespacing*{\section}{0pt}{1.4ex plus .3ex}{0.8ex}

\titleformat{\subsection}
  {\normalfont\small\bfseries\itshape}{\thesubsection.}{0.5em}{}
\titlespacing*{\subsection}{0pt}{1.2ex plus .2ex}{0.6ex}

%------------------------------------------------------------------------------
% Gaya caption
%------------------------------------------------------------------------------
\captionsetup{font=small,labelsep=period,justification=centering}
\captionsetup[table]{font={footnotesize,sc},skip=3pt}
\captionsetup[figure]{font=small,skip=4pt}

%------------------------------------------------------------------------------
% Header banner + footer
%------------------------------------------------------------------------------
\setlength{\headheight}{26pt}
\pagestyle{fancy}
\fancyhf{}
\renewcommand{\headrulewidth}{0.8pt}
\renewcommand{\footrulewidth}{0pt}
\fancyhead[L]{%
  \footnotesize\bfseries BULETIN PAGELARAN MAHASISWA NASIONAL BIDANG TEKNOLOGI INFORMASI DAN KOMUNIKASI\\%
  \mdseries\scriptsize ISSN: xxxxxx}
\fancyfoot[L]{\footnotesize Ketua Tim: Empat kata pertama judul makalah diikuti titik-titik (\ldots)}
\fancyfoot[R]{\footnotesize Volume 1 Nomor 1 September 2024}

\setlength{\parindent}{1.2em}
\setlength{\parskip}{0pt}

\begin{document}

%==============================================================================
% JUDUL, PENULIS, INTISARI (lebar penuh)
%==============================================================================
\twocolumn[{%
\begin{@twocolumnfalse}
\thispagestyle{fancy}
\vspace{0.5em}
\begin{center}
  {\LARGE\bfseries Peningkatan Keandalan dan Efisiensi Token LLM pada Data
   Time Series melalui Pendekatan Verifikasi Deterministik\par}
  \vspace{1.2em}
  {\normalsize Ahmad Naufal Farras\textsuperscript{1}, Jeremy Mattathias
   Mboe\textsuperscript{2}, Mochammad Henry Alifian\textsuperscript{3},
   Dhany Arifianto\textsuperscript{4}\par}
  \vspace{0.6em}
  {\footnotesize
   \textsuperscript{1}Fakultas Teknologi Elektro dan Informatika Cerdas,
   Institut Teknologi Sepuluh Nopember, Surabaya, 60111,
   email: 5054241038@student.its.ac.id\\
   \textsuperscript{2}Fakultas Teknologi Elektro dan Informatika Cerdas,
   Institut Teknologi Sepuluh Nopember, Surabaya, 60111,
   email: 5054241021@student.its.ac.id\\
   \textsuperscript{3}Fakultas Teknologi Elektro dan Informatika Cerdas,
   Institut Teknologi Sepuluh Nopember, Surabaya, 60111,
   email: 5054241024@student.its.ac.id\\
   \textsuperscript{4}Fakultas Kedokteran dan Kesehatan,
   Institut Teknologi Sepuluh Nopember, Surabaya, 60111\\[0.3em]
   \textit{Corresponding Author}: Ahmad Naufal Farras\par}
\end{center}
\vspace{1.0em}

{\small\justifying
\noindent\textbf{\textit{INTISARI}} --- Model bahasa besar atau LLM unggul
mengolah bahasa, tetapi lemah menalar atas angka. Karena itu model sering
keliru ketika menganalisis data \textit{time series}, misalnya menyebut tren,
lokasi lonjakan, atau besaran kenaikan yang tidak sesuai data, dengan kalimat
yang terdengar meyakinkan. Kekeliruan seperti ini berisiko pada bidang
bertaruh tinggi seperti kesehatan dan pangan. Penelitian ini mengusulkan
sebuah pendekatan penalaran \textit{time series} yang \textit{grounded},
artinya tiap langkah berpijak pada perhitungan atas deret asli, dan sekaligus
hemat \textit{token}. Ada dua kontribusi. Pertama, sebuah \textit{dataset}
penalaran \textit{time series} terverifikasi yang dibangun dari sumber publik
nyata di Indonesia, yaitu harga pangan PIHPS Bank Indonesia, kasus penyakit
Kementerian Kesehatan, cuaca BMKG, dan beban energi, lalu disimpan dalam
format JSONL dengan lapisan penalaran yang tiap angkanya sudah dicek. Kedua,
sebuah metode yang memasangkan LLM hasil \textit{fine-tune} kelas Qwen2.5-7B
dengan verifikator deterministik berbasis NumPy di luar model. LLM menuliskan
penalaran sebagai urutan operasi dari pustaka tetap, bukan prosa bebas, lalu
verifikator menghitung ulang tiap operasi atas deret asli dan membandingkannya
dengan angka yang diklaim memakai toleransi gabungan absolut dan relatif.
Metrik utamanya adalah skor \textit{grounding}, yaitu persentase langkah
penalaran yang angkanya tetap cocok saat dihitung ulang, dan dihitung tanpa
LLM penilai. Pada uji awal, verifikator berhasil mereproduksi penalaran jujur
atas deret kasus DBD mingguan dengan skor \textit{grounding} 100 persen, dan
menurunkan skor menjadi 66,7 persen ketika satu langkah dibuat berhalusinasi,
yaitu mengeklaim kenaikan 30 persen padahal 105 persen. Fondasi pustaka
operasi, skema, dan verifikator sudah lolos 20 pengujian. Pembeda penelitian
ini adalah penggabungan penalaran terverifikasi dan hemat \textit{token} pada
data \textit{time series} yang dibangun sendiri.
\par\vspace{0.7em}
\noindent\textbf{KATA KUNCI:} Penalaran \textit{Time Series}, Skor
\textit{Grounding}, Verifikator Deterministik, \textit{Chain-of-Thought},
Efisiensi \textit{Token}, \textit{Fine-tuning} LLM, QLoRA, Penalaran
Neuro-Simbolik, \textit{Data Mining}.
\par}
\vspace{1.5em}
\end{@twocolumnfalse}
}]

%==============================================================================
% ISI (dua kolom)
%==============================================================================
\justifying

\section{Pendahuluan}
Data \textit{time series} muncul di banyak keputusan penting, misalnya kasus
penyakit tiap minggu, harga pangan tiap hari, dan beban listrik tiap jam.
Selama ini pemakaian datanya lebih banyak diarahkan ke peramalan, yaitu
menebak angka berikutnya. Padahal angka ramalan saja sering belum cukup untuk
mengambil keputusan. Yang lebih dibutuhkan adalah tafsiran atas datanya,
umpamanya menilai apakah kenaikan pada beberapa minggu terakhir menandakan
wabah yang mulai terbentuk. Kemampuan menafsir semacam ini disebut penalaran,
dan porsinya masih tertinggal dibanding peramalan~\cite{merrill2024}.

Model bahasa besar dipakai sebagai antarmuka, tetapi lemah saat berhadapan
dengan angka. Pada deret panjang, model cenderung menebak dari pola bahasa
alih-alih menghitung, sehingga hasilnya bisa salah walau terdengar
meyakinkan~\cite{merrill2024,mtbench2025}. Sebagai contoh, untuk deret kasus
mingguan 12, 15, 11, 14, 13, 19, 25, 31, 29, 38, 45, 52, 61, 78 dengan
pertanyaan apakah ada indikasi wabah, jawaban yang benar adalah ada. Nilai
dasar pada awal deret sekitar 13 kasus, empat minggu terakhir naik dari 38 ke
78 atau sekitar 105 persen, dan nilai akhirnya kira-kira enam kali nilai
dasar. Pada kasus seperti ini model kerap keliru dalam tiga bentuk. Pertama,
menutupi bahaya dengan menyimpulkan kasus stabil padahal sudah naik ke 78.
Kedua, salah menunjuk lokasi lonjakan, umpamanya menyebut minggu semua padahal
nilainya baru naik minggu 18. Ketiga, salah menaksir besaran, umpamanya
menyebut kenaikan sekitar 30 persen padahal mencapai ratusan persen.

Kekeliruan seperti itu berisiko karena kalimatnya tampak rapi. Pada bidang
bertaruh tinggi seperti kesehatan dan pangan, jawaban tanpa alasan yang bisa
diperiksa tidak layak dipercaya. Karena itu penalaran perlu bersifat
\textit{grounded}, yaitu tiap langkahnya berpijak pada perhitungan atas deret
dan bisa diperiksa ulang, sehingga orang dapat menelusuri jalan pikirannya
sebelum bertindak. Pendekatan ini juga menyatukan dua hal yang selama ini
terpisah, yaitu metode statistik yang bisa menghitung tetapi tidak
menjelaskan, dan LLM yang bisa menjelaskan tetapi tidak menghitung.

\subsection{Kesenjangan Penelitian}
Penalaran atas \textit{time series} sedang ramai diteliti pada 2026, terlihat
dari sejumlah naskah di ICLR 2026, revisi tolok ukur, dan survei
khusus~\cite{mtbench2025,timeomni2026}. Ada pula temuan bahwa cara penalaran
yang berhasil di soal matematika belum tentu bekerja pada deret
waktu~\cite{merrill2024}. Arah penalaran yang lebih diverifikasi untuk
\textit{time series} memang sudah mulai dikerjakan~\cite{veritime2026}, tetapi
penalaran yang dibiarkan panjang dan boros \textit{token}~\cite{selfbudgeter2025}.
Di sisi lain, riset penghematan \textit{token} lewat pengaturan panjang
penalaran menurut tingkat kesulitan sudah berkembang, tetapi masih terbatas
pada soal matematika~\cite{selfbudgeter2025}. Sampai sekarang belum ada yang
menyatukan keduanya untuk \textit{time series}, dan belum ada \textit{dataset}
penalaran \textit{time series} terverifikasi yang dibangun dari sumber lokal.
Kami tidak mengeklaim bidang ini belum tersentuh, tetapi penggabungan penalaran
terverifikasi dan hemat \textit{token} di atas data lokal buatan sendiri memang
belum pernah dikerjakan.

\subsection{Kontribusi}
Penelitian ini memberi tiga kontribusi. Pertama, sebuah \textit{dataset}, yaitu
korpus penalaran \textit{time series} terverifikasi dari sumber publik nyata di
Indonesia, dalam format JSONL, dengan penalaran yang tiap langkahnya sudah dicek
ke angka. Bagian ini menjadi kontribusi utama karena usaha membangun datanya
terlihat jelas. Kedua, sebuah metode, yaitu pemasangan LLM hasil
\textit{fine-tune} dengan verifikator deterministik di luar model, ditambah
pengaturan panjang penalaran yang adaptif untuk menghemat \textit{token}.
Ketiga, sebuah metrik, yaitu skor \textit{grounding} yang mengukur kejujuran
penalaran terhadap angka dan dihitung tanpa LLM penilai. Kebaruannya bukan LLM,
sebab LLM meminjam yang sudah ada lalu disetel ulang, melainkan pada data yang
dibangun sendiri dan pada pasangan model dan verifikator yang membuat penalaran
bisa diperiksa sekaligus hemat.

\subsection{Rumusan Masalah}
Empat pertanyaan penelitian diajukan. RQ1, seberapa sering model sekarang
berhalusinasi saat menalar \textit{time series}, terlihat dari skor
\textit{grounding} rendah meski jawaban kadang benar. RQ2, apakah metode kami
menaikkan skor \textit{grounding} tanpa menurunkan akurasi jawaban. RQ3, apakah
penalaran adaptif kami lebih hemat \textit{token} dibanding penalaran panjang
pada akurasi setara. RQ4, komponen mana yang paling berpengaruh pada hasil.

\section{Tinjauan Pustaka}

\subsection{Penalaran atas Time Series}
Model bahasa besar masih sulit menalar tanpa contoh tentang \textit{time
series}, dan penalaran yang panjang tidak selalu menaikkan mutu analisis deret
waktu~\cite{merrill2024}. Tolok ukur multimodal untuk penalaran temporal
memperlihatkan model masih gagal ketika harus menggabungkan angka dan
teks~\cite{mtbench2025}. Upaya ini mendorong penalaran yang lebih rumit atas
\textit{time series} lewat pelatihan berbasis insentif~\cite{timeomni2026}.
Kumpulan karya ini menunjukkan bidangnya aktif tetapi belum tuntas.

\subsection{Chain-of-Thought dan Program-of-Thought}
Memunculkan langkah penalaran antara terbukti menaikkan hasil pada tugas rumit,
tetapi langkah itu tetap berupa prosa yang bisa salah hitung~\cite{cot2022}.
Pemindahan komputasi ke penafsir program membuat jawaban dihitung alih-alih
ditebak~\cite{pal2023}. Penelitian ini memakai semangat yang sama, yaitu LLM
tidak pernah memegang angka, ia hanya mengusulkan operasi, dan program di luar
model yang menghitung dan mengecek.

\subsection{Penalaran Terverifikasi untuk Time Series}
Arah yang paling dekat dengan penelitian ini menggarap penalaran \textit{time
series} yang dapat diverifikasi melalui sintesis dan penjadwalan data berpikir,
yang dikenal VeriTime~\cite{veritime2026}. Penelitian kami diletakkan langsung
terhadap arah tersebut, dengan pembeda menggabungkan penalaran terverifikasi dan
hemat \textit{token} serta membangun data sendiri dari konteks lokal. Klaim
bahwa arah terverifikasi cenderung menghasilkan penalaran panjang dan mahal
perlu dicek terhadap naskah aslinya sebelum dipakai sebagai poros pembeda.

\subsection{Efisiensi Token dan Penalaran Adaptif}
Panjang penalaran dapat dibuat hemat dan menyesuaikan tingkat kesulitan lewat
pengaturan anggaran \textit{token}, meski sejauh ini terbatas pada soal
matematika~\cite{selfbudgeter2025}. Penalaran panjang berbasis pembelajaran
penguatan juga menjadi lazim setelah 2025~\cite{deepseekr1}. Gagasan anggaran
adaptif itu kami pinjam dan bawa ke ranah \textit{time series}, yang belum
pernah dilakukan.

\subsection{Pelatihan Hemat Sumber Daya}
Adaptasi peringkat rendah dan kuantisasi NF4 4-bit memungkinkan penyetelan model
7B pada GPU kelas T4 atau V100~\cite{lora2022,qlora2023}. Untuk deret panjang,
pendekatan penyandi bergaya \textit{reprogramming} dapat menjadi opsi
tambahan~\cite{timellm2024}. Karena kekuatan proyek ada pada data dan
\textit{grounding}, model kelas 7B seperti Qwen2.5 dinilai memadai tanpa perlu
ukuran besar~\cite{qwen2024}.

\section{Usulan Penelitian}
Penelitian ini memisahkan tiga hal yang jangan tertukar, yaitu data sebagai
bahan, metode sebagai cara membuat penalaran yang \textit{grounded} dan hemat,
serta metrik sebagai alat ukur. Metodenya terdiri atas dua komponen berbeda
sifat, yaitu komponen model berupa LLM yang disetel dan komponen berupa
verifikator deterministik yang tidak dilatih. Posisi kami terhadap karya terkait
diringkas pada Tabel~\ref{tab:posisi}. Celah yang kami ambil adalah penggabungan
penalaran terverifikasi dan sekaligus hemat \textit{token} di atas data lokal.

\begin{table}[t]
\centering
\caption{Posisi Penelitian Ini terhadap Karya Terkait}
\label{tab:posisi}
\footnotesize
\begin{tabular}{@{}lcccc@{}}
\toprule
\textbf{Karya} & \textbf{Terverifikasi} & \textbf{Hemat} & \textbf{\textit{Time}} & \textbf{Data} \\
 & & \textbf{\textit{token}} & \textbf{\textit{series}} & \textbf{lokal} \\
\midrule
CoT / PoT \cite{cot2022,pal2023} & sebagian & tidak & tidak & -- \\
SelfBudgeter \cite{selfbudgeter2025} & tidak & ya & tidak & -- \\
VeriTime \cite{veritime2026} & ya & tidak & ya & tidak \\
Usulan kami & ya & ya & ya & ya \\
\bottomrule
\end{tabular}
\end{table}

\section{Metodologi}

\subsection{Gambaran Sistem}
Alur inti sistem berjalan sebagai berikut. LLM menerima deret angka beserta
pertanyaan, lalu menuliskan penalaran dalam bentuk langkah beroperasi, bukan
prosa bebas. Verifikator deterministik di luar model kemudian menghitung ulang
tiap operasi atas deret asli dan mengecek angka yang diklaim cocok dalam
toleransi. LLM yang menalar, program yang mengoreksi. Keduanya benda terpisah
dan bekerja bersama, bukan program yang ditanam ke dalam LLM. Gambaran alur ini
ada pada Gambar~\ref{fig:arsitektur}.

\begin{figure}[t]
\centering
\includegraphics[width=0.92\linewidth]{arsitektur_sistem.png}
\caption{Arsitektur sistem: pemisahan komponen model (LLM penalar) dan komponen
program (verifikator deterministik).}
\label{fig:arsitektur}
\end{figure}

\subsection{Konstruksi Dataset}
Deret angka diambil dari sumber publik nyata lewat \textit{scraping} atau API,
jadi bukan \textit{dataset} siap pakai (Tabel~\ref{tab:sumber}). Untuk tiap
deret, kami mengajukan pertanyaan, menghitung jawaban benar dengan rumus, lalu
menyusun langkah penalaran yang tiap angkanya dicek ke deret. Langkah dibatasi
pada pustaka operasi yang sudah ditetapkan agar tiap langkah bisa dihitung
ulang. Untuk mencapai skala, sebagian besar data latih dibuat semi-sintetis,
yaitu deret nyata atau deret terkendali yang dipasangkan dengan penalaran yang
diperiksa otomatis, sedangkan data uji didominasi deret asli agar terbukti
bekerja di lapangan. Pembagian dan distratifikasi menurut jenis tugas dan
tingkat kesulitan, dan sumber pada data uji dijaga tidak muncul pada data latih
supaya tidak bocor.

\begin{table}[t]
\centering
\caption{Sumber Deret Publik yang Direncanakan}
\label{tab:sumber}
\footnotesize
\begin{tabular}{@{}p{1.3cm}p{1.6cm}p{1.5cm}p{1.1cm}p{1.6cm}@{}}
\toprule
\textbf{Domain} & \textbf{Sumber} & \textbf{Metode akses} & \textbf{Frekuensi} & \textbf{Contoh variabel} \\
\midrule
Harga pangan & PIHPS Bank Indonesia & Scraping / API & Harian & Beras, cabai, minyak goreng \\
Kesehatan & Dashboard Kemenkes & Scraping & Mingguan & Kasus DBD, ISPA, diare \\
Cuaca & BMKG & API / scraping & Harian, jam & Suhu, curah hujan, kelembapan \\
Energi & Sumber beban terbuka & API & Jam & Beban dan konsumsi listrik \\
\bottomrule
\end{tabular}
\end{table}

\subsection{Skema JSONL dan Pustaka Operasi}
Data disimpan sebagai JSONL, satu sampel per baris. Tiap sampel memuat empat
bagian, yaitu deret berisi nama, satuan, frekuensi, dan konteks; pertanyaan;
daftar penalaran yang tiap langkahnya berisi operasi, hasil yang diklaim, dan
teks; serta jawaban berisi label dan tingkat keyakinan. Tiap langkah menamai
satu operasi dari pustaka tetap pada Tabel~\ref{tab:operasi} agar bisa
dijalankan ulang oleh verifikator. Operasi majemuk dapat disusun dari operasi
dasar.

\begin{table}[t]
\centering
\caption{Pustaka Operasi Dasar}
\label{tab:operasi}
\footnotesize
\begin{tabular}{@{}>{\RaggedRight\arraybackslash}p{2.1cm}>{\RaggedRight\arraybackslash}p{2.2cm}>{\RaggedRight\arraybackslash}p{2.9cm}@{}}
\toprule
\textbf{Operasi} & \textbf{Arti} & \textbf{Contoh} \\
\midrule
rata2(rentang) & rata-rata atas rentang & rata2(nilai[0:5]) \\
delta(a, b) & selisih dua titik & delta(nilai[0], nilai[15]) \\
persen\_naik(a, b) & perubahan persen & persen\_naik(nilai[11], nilai[15]) \\
rasio(a, b) & perbandingan dua nilai & rasio(nilai[15], baseline) \\
slope(rentang) & kemiringan tren linear & slope(nilai[10:16]) \\
min / max(rentang) & nilai ekstrem & max(nilai) \\
z\_score(titik) & deviasi dalam simpangan baku & z\_score(nilai[15]) \\
deteksi\_\allowbreak anomali(ambang) & titik melewati ambang & deteksi\_\allowbreak anomali(z=3) \\
bandingkan\_\allowbreak segmen(r1, r2) & banding dua segmen & bandingkan\_\allowbreak segmen(\allowbreak nilai[0:8],\allowbreak{} nilai[8:15]) \\
\bottomrule
\end{tabular}
\end{table}

\subsection{Verifikator Deterministik dan Skor Grounding}
Verifikator adalah program di luar LLM, ditulis dengan Python dan NumPy dan
tidak dilatih. Program ini membaca \textit{string} operasi tiap langkah,
menerjemahkan tokennya (nilai[a:b] menjadi irisan, nilai[i] menjadi satu elemen,
angka menjadi literal, dan nama menjadi hasil langkah sebelumnya), menghitung
ulang atas deret asli, lalu membandingkan hasilnya dengan angka yang diklaim.
Satu langkah disebut \textit{grounded} bila memenuhi persamaan~\eqref{eq:grounded}.

\begin{equation}
\label{eq:grounded}
\left| v_{\text{hitung}} - v_{\text{klaim}} \right|
\le \max\!\left( \varepsilon_{\text{abs}},\;
\varepsilon_{\text{rel}} \cdot \left| v_{\text{hitung}} \right| \right)
\end{equation}

dengan nilai bawaan kedua toleransi, absolut dan relatif, sama dengan 0,01. Suku
relatif berperan sebagai penyetel kalibrasi. Suku ini membuat nilai tersaji
105,3 tetap lolos terhadap hasil hitung ulang 105,26 yang hanya beda pembulatan,
sementara salah besaran yang nyata, misalnya 30 lawan 105, tetap gagal. Skor
\textit{grounding} atas sampel dihitung dengan persamaan~\eqref{eq:skor}.

\begin{equation}
\label{eq:skor}
\text{Skor grounding} = \frac{n_{\text{grounded}}}{n_{\text{skalar}}} \times 100\%
\end{equation}

Langkah bukan skalar, misalnya deteksi\_anomali yang mengembalikan daftar indeks,
untuk sementara tidak dihitung dalam skor skalar. Verifikator dipakai di tiga
tempat, yaitu membersihkan label saat membangun data dengan membuang langkah
yang gagal, menghitung skor \textit{grounding} saat evaluasi, dan secara
opsional menjadi sinyal imbalan saat pelatihan.

\subsection{Penalaran Adaptif}
Panjang penalaran diatur menurut kesulitan. Deret dengan pola jelas cukup diberi
penalaran pendek, sedangkan deret yang ambigu diberi penalaran panjang. Cara ini
menekan jumlah \textit{token} tanpa menurunkan akurasi, mengikuti gagasan
anggaran \textit{token} adaptif~\cite{selfbudgeter2025} yang kami bawa ke ranah
\textit{time series}.

\subsection{Protokol Pelatihan}
Penyetelan memakai QLoRA 4-bit agar muat pada GPU T4 atau V100
(Tabel~\ref{tab:latih}). Model kelas 7B dinilai cukup karena kekuatan proyek ada
pada data dan \textit{grounding}. Sebagai perluasan opsional, pembelajaran
penguatan dengan imbalan terverifikasi memberi nilai lebih bila langkah
penalaran lolos pemeriksaan \textit{grounding}.

\begin{table}[t]
\centering
\caption{Konfigurasi Pelatihan Kandidat}
\label{tab:latih}
\footnotesize
\begin{tabular}{@{}p{1.8cm}p{2.0cm}p{1.4cm}p{1.6cm}@{}}
\toprule
\textbf{Parameter} & \textbf{Nilai} & \textbf{Parameter} & \textbf{Nilai} \\
\midrule
Metode & QLoRA (NF4 4-bit) & Penjadwal & cosine + warmup \\
Peringkat LoRA & 16 sampai 64 & Presisi & bf16 \\
Alfa LoRA & 32 & Epoch & 2 sampai 3 \\
Laju belajar & 1e-4 sampai 2e-4 & Kerangka & LLaMA-Factory, ZeRO-2 \\
\bottomrule
\end{tabular}
\end{table}

\section{Hasil Awal dan Rancangan Eksperimen}
Catatan penting. Penelitian masih pada Fase 1 atau fondasi. Bagian V.A sampai
V.C melaporkan bukti kecil nyata dan bisa direproduksi dari verifikator,
sedangkan bagian V.D dan V.E memaparkan rancangan eksperimen penuh yang akan
dijalankan pada Fase 4 dan 5. Tidak ada angka hasil model terlatih yang diklaim
pada draf ini.

\subsection{Implementasi Fondasi}
Empat modul inti sudah dibuat dan diuji, yaitu skema JSONL, pustaka operasi
berbasis NumPy, verifikator, dan modul metrik. Seluruh fondasi lolos 20
pengujian otomatis yang mencakup kebenaran tiap operasi, baca tulis JSONL
bolak-balik, validasi skema, serta perhitungan skor \textit{grounding}.

\subsection{Studi Kasus: Deteksi Wabah DBD}
Untuk deret kasus DBD mingguan sebuah kabupaten selama 16 minggu pada
Tabel~\ref{tab:dbd100}, penalaran jujur menyusun tiga langkah. Verifikator
menghitung ulang tiap operasi atas deret asli lalu mencocokkannya dengan angka
yang diklaim. Ketiganya \textit{grounded}, sehingga skor \textit{grounding}
mencapai 100 persen dan menurunkan kesimpulan wabah dengan keyakinan tinggi.

\begin{table}[t]
\centering
\caption{Reproduksi Penalaran Jujur, Skor Grounding 100 Persen}
\label{tab:dbd100}
\footnotesize
\begin{tabular}{@{}cp{3.0cm}ccc@{}}
\toprule
\textbf{Lgk.} & \textbf{Operasi} & \textbf{Diklaim} & \textbf{Hitung ulang} & \textbf{Grounded} \\
\midrule
1 & rata2(nilai[0:5]) & 13,0 & 13,0 & ya \\
2 & persen\_naik(nilai[11], nilai[15]) & 105,3 & 105,26 & ya \\
3 & rasio(nilai[15], 13,0) & 6,0 & 6,0 & ya \\
\bottomrule
\end{tabular}
\end{table}

\subsection{Validasi Deteksi Halusinasi}
Untuk menguji apakah metrik menangkap kesalahan nyata, langkah kedua diubah
menjadi klaim salah besaran, yaitu 30 persen padahal 105 persen. Verifikator
menandai langkah itu tidak \textit{grounded} karena selisih 75,26 jauh
melampaui toleransi pada persamaan~\eqref{eq:grounded}, sehingga skor
\textit{grounding} turun menjadi 66,7 persen (Tabel~\ref{tab:dbd66}). Inilah
metrik yang menjadi tulang punggung proyek, yaitu kejujuran penalaran terhadap
angka yang dihitung tanpa LLM penilai.

\begin{table}[t]
\centering
\caption{Deteksi Halusinasi Salah Besaran, Skor Grounding 66,7 Persen}
\label{tab:dbd66}
\footnotesize
\begin{tabular}{@{}cp{3.0cm}ccc@{}}
\toprule
\textbf{Lgk.} & \textbf{Operasi} & \textbf{Diklaim} & \textbf{Hitung ulang} & \textbf{Grounded} \\
\midrule
1 & rata2(nilai[0:5]) & 13,0 & 13,0 & ya \\
2 & persen\_naik(nilai[11], nilai[15]) & 30,0 & 105,26 & tidak \\
3 & rasio(nilai[15], 13,0) & 6,0 & 6,0 & ya \\
\bottomrule
\end{tabular}
\end{table}

\subsection{Rancangan Eksperimen Penuh}
Empat eksperimen dirancang untuk menjawab RQ1 sampai RQ4. Pertama, uji
halusinasi dasar dengan menguji model tanpa penyetelan pada tolok ukur kami
untuk membuktikan skor \textit{grounding} rendah. Kedua, uji metode lawan
pembanding untuk menunjukkan skor \textit{grounding} naik tanpa menurunkan
akurasi. Ketiga, uji efisiensi dengan membandingkan \textit{token} penalaran
adaptif lawan penalaran panjang pada akurasi setara. Keempat, uji ablasi dengan
mematikan komponen satu per satu, yaitu verifikator saat pelatihan, penalaran
adaptif, dan format beroperasi, untuk mengukur sumbangan tiap bagian.
Pembandingnya meliputi LLM umum tanpa penyetelan yang menulis prosa bebas, LLM
dengan penalaran panjang bergaya terverifikasi wakil arah
VeriTime~\cite{veritime2026}, dan metode statistik murni sebagai batas atas
hitungan tanpa penjelasan.

\subsection{Metrik dan Hasil yang Diharapkan}
Tiga metrik dipakai dan semuanya deterministik, yaitu akurasi jawaban berupa
persentase label akhir yang benar, skor \textit{grounding} pada
persamaan~\eqref{eq:skor} sebagai metrik utama, dan efisiensi \textit{token}
berupa rata-rata \textit{token} per sampel yang dilaporkan bersama akurasi agar
imbangnya terlihat. Hasil yang diharapkan meliputi \textit{dataset} penalaran
\textit{time series} terverifikasi yang siap dipakai ulang, bukti bahwa model
sekarang berskor \textit{grounding} rendah, bukti bahwa metode kami menaikkan
\textit{grounding} tanpa menurunkan akurasi dan lebih hemat \textit{token},
serta prototipe utuh yang bisa menjawab pertanyaan atas deret dengan penalaran
yang bisa diperiksa.

\section{Kesimpulan dan Saran}

\subsection{Kesimpulan}
Draf ini menetapkan fondasi sebuah pendekatan penalaran \textit{time series}
yang memisahkan LLM penalar dari verifikator deterministik, dengan penalaran
berbentuk operasi yang tiap angkanya bisa dihitung ulang. Fondasinya sudah
terbukti pada studi kasus nyata. Verifikator mereproduksi penalaran jujur pada
deret DBD dengan skor \textit{grounding} 100 persen dan menandai penalaran
berhalusinasi dengan skor turun ke 66,7 persen, serta lolos 20 pengujian.
Sumbangan untuk kompetisi data mining adalah data lokal terverifikasi yang usaha
pembuatannya besar, dipasangkan dengan metode yang bisa diaudit dan hemat
\textit{token}.

\subsection{Saran dan Rencana Kerja}
Pekerjaan berikutnya mengikuti enam fase, yaitu finalisasi pustaka operasi dan
skema pada Fase 1, penarikan dari sumber publik pada Fase 2, sintesis dan
verifikasi lapisan penalaran pada Fase 3, penyetelan model dan integrasi alur
pada Fase 4, eksperimen serta ablasi dan evaluasi pada Fase 5, dan pembuatan
prototipe serta penulisan laporan pada Fase 6. Titik paling menentukan ada di
Fase 1 dan 3 karena pustaka operasi dan verifikasi label menentukan seluruh
data. Risiko utama berupa tumpang tindih dengan VeriTime ditekan dengan
menegaskan gabungan penalaran terverifikasi dan hemat \textit{token} serta data
lokal buatan sendiri. Dari sisi etika, seluruh sumber bersifat publik, sistem
diposisikan sebagai penunjang keputusan dan bukan nasihat medis atau finansial
pribadi, dan penalaran yang \textit{grounded} justru mendukung
pertanggungjawaban karena jejak hitungnya bisa ditelusuri.

%==============================================================================
% REFERENSI
%==============================================================================
\begin{thebibliography}{99}
\footnotesize
\bibitem{merrill2024}
M. Merrill, M. Tan, V. Gupta, T. Hartvigsen, dan T. Althoff, ``Language Models
Still Struggle to Zero-shot Reason about Time Series,'' dalam \textit{Findings
of EMNLP}, 2024.

\bibitem{mtbench2025}
``MTBench: A Multimodal Time Series Benchmark for Temporal Reasoning and
Question Answering,'' arXiv:2503.16858, 2025.

\bibitem{veritime2026}
``VeriTime: Time Series Reasoning via Process-Verifiable Thinking Data Synthesis
and Scheduling for Tailored LLM Reasoning,'' arXiv:2602.07830, 2026.

\bibitem{selfbudgeter2025}
``SelfBudgeter: Adaptive Token Allocation for Efficient LLM Reasoning,''
arXiv:2505.11274, 2025.

\bibitem{timeomni2026}
``TimeOmni-1: Incentivizing Complex Reasoning with Time Series in Large Language
Models,'' dalam \textit{Proc. ICLR}, 2026.

\bibitem{cot2022}
J. Wei et al., ``Chain-of-Thought Prompting Elicits Reasoning in Large Language
Models,'' dalam \textit{Proc. NeurIPS}, 2022.

\bibitem{pal2023}
L. Gao et al., ``PAL: Program-aided Language Models,'' dalam \textit{Proc.
ICML}, 2023.

\bibitem{lora2022}
E. J. Hu et al., ``LoRA: Low-Rank Adaptation of Large Language Models,'' dalam
\textit{Proc. ICLR}, 2022.

\bibitem{qlora2023}
T. Dettmers, A. Pagnoni, A. Holtzman, dan L. Zettlemoyer, ``QLoRA: Efficient
Finetuning of Quantized LLMs,'' dalam \textit{Proc. NeurIPS}, 2023.

\bibitem{deepseekr1}
DeepSeek-AI, ``DeepSeek-R1: Incentivizing Reasoning Capability in LLMs via
Reinforcement Learning,'' arXiv:2501.12948, 2025.

\bibitem{qwen2024}
Qwen Team, ``Qwen2.5 Technical Report,'' arXiv:2412.15115, 2024.

\bibitem{timellm2024}
M. Jin et al., ``Time-LLM: Time Series Forecasting by Reprogramming Large
Language Models,'' dalam \textit{Proc. ICLR}, 2024.

\bibitem{llamafactory2024}
Y. Zheng et al., ``LLaMA-Factory: Unified Efficient Fine-Tuning of 100+ Language
Models,'' dalam \textit{Proc. ACL System Demonstrations}, 2024.

\bibitem{bi2024}
Bank Indonesia, ``Pusat Informasi Harga Pangan Strategis Nasional,'' tersedia di
\url{https://www.bi.go.id/hargapangan}, diakses [tanggal].

\bibitem{kemenkes}
Kementerian Kesehatan RI, ``Dashboard Data Penyakit,'' tersedia di [URL],
diakses [tanggal].

\bibitem{bmkg}
BMKG, ``Data Online Pusat Database BMKG,'' tersedia di
\url{https://dataonline.bmkg.go.id}, diakses [tanggal].
\end{thebibliography}

\end{document}
